% Part III: ECGs and imaging as data. Sources: docs/importing.md "Files"; docs/administration.md
% "Files and recordings", "Import formats"; viewers/ecg/python/muse/README.md;
% viewers/pacs/python/dicom/README.md; docs/edf.md; docs/security.md 7.3 (the file rules) and
% docs/protecting-patient-data.md#dicom (the DICOM element rules). Checked against
% medorama/src/tables.rs (files of rows), patient/files.rs (kinds, viewers, hints), gate.rs
% (FileKind), viewers/ecg/python/muse/muse_reader_py/cli.py, viewers/pacs/python/dicom.
\chapter{ECGs and imaging as data}
\label{chap:ecg-imaging}

This chapter explains how ECGs, images and documents come into Medorama, and how they are
opened and protected. It is written for data managers who build the tables that name these
files.

Resting ECGs, echocardiograms, angiograms, CT and MR series, X-rays, reports and documents come
into Medorama as the files of a table's rows (\cref{chap:importing}). Each row names its
patient, its time and its file, and becomes an item of the ECG, Echo, Cath lab or Imaging column
of the Patient Journey. Its file opens in the ECG viewer (\cref{chap:ecg-viewer}) or the PACS
viewer (\cref{chap:pacs-viewer}). Such tables are written by the systems that hold the files, or
made from the files themselves with \cmd{muse-py index} and \cmd{dicom-py index}. Waveform
recordings (ICU monitoring, Holters, EEGs and sleep studies) are the subject of
\cref{chap:recordings}.

\section{Files as attachments}
\label{sec:ecg-imaging-files}

A row's file is the \emph{attachment} of the row's item when the row's category (or the table's)
is not a recording category (\code{icu}, \code{holter}, \code{eeg} or \code{sleep}). Its kind is
recognized from its extension, and a file without one is recognized as DICOM by the \code{DICM}
marker after its 128-byte preamble (\cref{tab:ecg-imaging-viewers}). The file must lie inside
\code{data\_roots}. It is read where it is, never copied, and shown only to users who may see the
patient.

\begin{table}[htbp]
\caption{Which viewer opens an attachment, by the item's category and the file's kind.}
\label{tab:ecg-imaging-viewers}
\begin{tabularx}{\linewidth}{@{}P{2.6cm}LP{2.6cm}@{}}
\toprule
{\tabhead Category} & {\tabhead Files} & {\tabhead Opened in} \\
\midrule
ECG & A WFDB header (\file{.hea}) with its signal files, GE MUSE XML (\file{.xml}), a DICOM ECG
  (\file{.dcm}, \file{.dicom} or no extension), an EDF or BDF file (\file{.edf}, \file{.bdf}), a
  JSON ECG (\file{.json}), a picture (\file{.png}, \file{.jpg}, \file{.jpeg}, \file{.gif},
  \file{.webp}, \file{.svg}), a PDF & The ECG viewer \\
Echo, Cath lab, Imaging & DICOM files, a folder (a DICOM series, or pictures that make one
  loop), videos (\file{.mp4}, \file{.webm}, \file{.mov}, \file{.m4v}), pictures, PDF & The PACS
  viewer \\
Any other & A picture, a PDF, a video, a text file (\file{.txt}) & The page itself \\
\bottomrule
\end{tabularx}
\end{table}

The item's name is the row's \code{name}, or else the file's name without its extension. Its
measurements are the row's measurement columns (\code{QTc (ms)}, \code{LVEF (\%)}): numbers that
cohorts search, in their units. The row's text columns (\code{interpretation}, \code{report},
\code{findings} \dots) are the item's text.

\section{ECGs}
\label{sec:ecg-imaging-ecg}

\screenshot{ecg.png}{A synthetic resting 12-lead ECG in the ECG viewer (light theme), opened from its item.}{fig:ecg-imaging-ecg}

The ECG viewer opens the following formats (\cref{chap:ecg-viewer} describes how it reads each):
\begin{itemize}
\item \textbf{WFDB}: a record's header, with the signal files it names, which the item's file
  brings with it;
\item \textbf{GE MUSE XML}: a resting ECG as the MUSE system stores it, with its rhythm waveform,
  measurements and statements;
\item \textbf{DICOM ECGs}: 12-lead and Holter waveforms;
\item \textbf{EDF, EDF+, BDF and BDF+}: the first 10 seconds of the ECG signals
  (\cref{sec:ecg-viewer-edf}); long recordings belong to a recording category and open in
  VitalStream instead;
\item \textbf{JSON}: the ECG viewer's own form (\code{name}, \code{fs}, \code{leads}, each with
  its \code{name}, \code{units} and \code{samples}, and \code{meta}), as \cmd{muse-py json} writes
  it;
\item \textbf{pictures and PDF}: an ECG printed or scanned, shown as it is, with its scale taken
  from the row's columns (\cref{sec:ecg-imaging-hints}).
\end{itemize}

A table of ECGs names, in each row, the patient, the time, the file and, when there are any, the
measurements and statements. The following commands import such a table:
\begin{lstlisting}[language=CSV]
patient_id,encounter_id,time,QT (ms),QTc (ms),interpretation,file
P900001,E9100002,2024-01-02 08:00,410,452,Sinus tachycardia,ecg/ecg_P900001_20240102.hea
\end{lstlisting}
\begin{lstlisting}[style=shell]
medorama import-table ecg_index.csv              # the category from the file name: ecg
medorama import-table ecgs.csv --category ecg    # else name it
\end{lstlisting}
An export of GE MUSE's ECG metadata (one row per ECG, the XML elements as columns) is recognized
and converted by the import format \code{ge-muse} (\code{ge-muse-patient-id} when its patients
are named by a research ID: \cref{chap:importing}).

\subsection{Building an ECG database: muse-py index}
\label{sec:ecg-imaging-muse}

A folder of GE MUSE XML files becomes an ECG database with \cmd{muse-py index}, the MUSE reader of
the ECG viewer, which needs no other package. The following commands install it, index a folder
and import the result:
\begin{lstlisting}[style=shell]
pip install -e viewers/ecg/python/muse
muse-py index /data/muse -o /data/muse/ecg.csv     # one row per MUSE file
medorama import-table /data/muse/ecg.csv           # each row an ECG of its patient
\end{lstlisting}
It reads every \file{.xml} file below the folder (hidden files left out), and writes one row per
ECG with the columns of \cref{tab:ecg-imaging-muse}; the files it cannot read are listed on
standard error. File paths are relative to the CSV, which the import reads from the CSV's folder.

\begin{table}[htbp]
\caption{The columns \cmd{muse-py index} writes.}
\label{tab:ecg-imaging-muse}
\begin{tabularx}{\linewidth}{@{}P{4.4cm}L@{}}
\toprule
{\tabhead Column} & {\tabhead Holds} \\
\midrule
\code{mrn} & The MUSE \code{PatientDemographics/PatientID}, in a column named \code{mrn}:
  Medorama links each row to its patient by it, as an MRN\phimark \\
\code{time} & The acquisition date and time, as the cart's clock wrote them \\
\code{category}, \code{name} & \code{ecg}, \code{12-lead ECG} \\
\code{value}, \code{unit} & The ventricular rate, \code{bpm} \\
\code{Vent. rate (bpm)}, \code{Atrial rate (bpm)}, \code{PR (ms)}, \code{QRS (ms)},
  \code{QT (ms)}, \code{QTc (ms)}, \code{QTcF (ms)}, \code{P axis (\textdegree)},
  \code{R axis (\textdegree)}, \code{T axis (\textdegree)}, \code{QRS count (beats)} & The
  measurements (\code{RestingECGMeasurements}) \\
\code{interpretation} & The computer's statements (Marquette 12SL) \\
\code{confirmed} & The physician's statements: confirmed, else edited \\
\code{file} & The MUSE file, relative to the CSV \\
\bottomrule
\end{tabularx}
\end{table}

\cmd{muse-py wfdb ECG.xml OUT\_DIR} converts one MUSE file into a WFDB record (format 16, 1\,$\mu$V
per unit), its details, measurements and statements in its header's comments; \cmd{dicom-py wfdb
ECG.dcm OUT\_DIR} does the same for a DICOM ECG (12-lead or Holter). Neither writes the patient's
identity into the record.

\section{Imaging}
\label{sec:ecg-imaging-imaging}

\screenshot{echo.png}{A synthetic echocardiogram loop in the PACS viewer, opened from its item.}{fig:ecg-imaging-echo}

The PACS viewer opens DICOM studies of every modality (echocardiography and other ultrasound,
angiography and fluoroscopy, IVUS and OCT, CT, MR, nuclear medicine, X-rays), DICOM structured
reports and encapsulated PDFs, videos, pictures and PDF documents (\cref{chap:pacs-viewer}). A
row's file may be:
\begin{itemize}
\item \textbf{a DICOM file}: a multi-frame loop or run, a single image, a structured report, an
  encapsulated PDF;
\item \textbf{a folder}: a DICOM series (its files, with or without an extension), or a folder of
  pictures, which the viewer plays as one loop at the row's frame rate;
\item \textbf{a video, a picture or a PDF}.
\end{itemize}
A folder's files are its own, not its subfolders'; only the files of the item's folder, or of the
file and its WFDB signal files, can be read through the item.

\subsection{Building an imaging database: dicom-py index}
\label{sec:ecg-imaging-dicom}

A folder of DICOM files becomes an imaging database with \cmd{dicom-py index}, which uses pydicom
and numpy. The following commands install it, index a folder and import the result:
\begin{lstlisting}[style=shell]
pip install -e viewers/pacs/python/dicom
dicom-py index /data/echo -o /data/echo/imaging.csv   # one row per loop, run, series, report
medorama import-table /data/echo/imaging.csv          # the rows reach their patients
\end{lstlisting}
It reads every DICOM file under the folder, headers only, never the pixels, and writes one row per
object a viewer shows (\cref{tab:ecg-imaging-dicom}). The objects it cannot link are listed on
standard error.

\begin{table}[htbp]
\caption{The rows \cmd{dicom-py index} writes.}
\label{tab:ecg-imaging-dicom}
\begin{tabularx}{\linewidth}{@{}P{4.4cm}L@{}}
\toprule
{\tabhead Object} & {\tabhead Row} \\
\midrule
An ultrasound loop, an angiography run, any multi-frame file & One per file, with its frames and
  frame rate \\
A series of single-frame images (CT, MR, NM) & One per series, linking the series' folder: only
  when that folder holds one patient's files (else it is listed as not linked, and no row is
  written) \\
A single image (an X-ray) & One per file \\
A structured report (echo, cath \dots) & One per file; each measurement a column
  \code{Name (unit)}, such as \code{Left Ventricular Ejection Fraction (\%)} (a repeated
  measurement keeps its first value, or its mean when one is given); its text in
  \code{interpretation} \\
An encapsulated PDF & One per file \\
A DICOM ECG (12-lead, Holter) & One per file; its measurements as the ECG viewer names them
  (\code{QT (ms)}, \code{QRS (ms)}, \code{Vent. rate (bpm)} \dots), its statements in
  \code{interpretation} \\
\bottomrule
\end{tabularx}
\end{table}

The columns are \code{mrn} (the patient ID\phimark), \code{time} (the device's clock),
\code{category} and \code{name}, then \code{modality}, \code{description}, \code{study\_uid},
\code{series\_uid}, \code{images}, \code{frames} and \code{fps}; the measurement columns follow,
then \code{interpretation} and \code{file} (relative to the CSV). The category is \code{echo} for
ultrasound, \code{cath} for angiography, fluoroscopy, IVUS and OCT, \code{ecg} or \code{holter}
for ECG waveforms, and \code{img} for everything else; a report or a PDF takes its category from
the words of its title (echo, cath, ECG).

Columns of different databases are never merged: an echo database's \code{LVEF (\%)} and a cath
database's are distinct, and a cohort's criterion names which one (\cref{chap:cohorts}).

\section{Hints for the viewers}
\label{sec:ecg-imaging-hints}

A picture or a video does not say its scale, and a folder of pictures does not say its frame
rate. Columns of the item's row say them: the first of its columns whose name (without its unit)
and unit are in \cref{tab:ecg-imaging-hints}, and whose value is a positive number, gives the
hint. Names are compared ignoring case, spaces and punctuation.

\begin{table}[htbp]
\caption{The columns that give the viewers a scale and a frame rate.}
\label{tab:ecg-imaging-hints}
\begin{tabularx}{\linewidth}{@{}P{2.0cm}LP{4.4cm}@{}}
\toprule
{\tabhead Hint} & {\tabhead Column names} & {\tabhead Units} \\
\midrule
Scale (mm per pixel) & \code{pixel spacing}, \code{pixel size}, \code{spacing},
  \code{mm per pixel}, \code{mm/pixel}, \code{resolution} & \code{mm}, \code{mm/px},
  \code{mm/pixel} or none; \code{cm}, \code{cm/px}, \code{cm/pixel} (times 10); \code{um},
  \code{\textmu m} (divided by 1000) \\
Frame rate & \code{fps}, \code{frame rate}, \code{frames per second} & \code{fps}, \code{Hz},
  \code{/s}, \code{1/s} or none \\
\bottomrule
\end{tabularx}
\end{table}

For example, \code{pixel spacing (mm)} and \code{frame rate (fps)}, or the \code{fps} column
\cmd{dicom-py index} writes.

\section{In Privacy Mode}
\label{sec:ecg-imaging-privacy}

\begin{phibox}
Files are protected health information: their headers, their names, and the text burned into
their pixels. In Privacy Mode, which always applies to users without the PHI permission, every
file leaves through the privacy gate. The gate de-identifies what it can, passes images with a
warning, and refuses the rest. The rules are given in full in \cref{sec:protecting-files}, and
the DICOM rules element by element in \cref{sec:protecting-dicom}.
\end{phibox}

\Cref{tab:ecg-imaging-gate} summarizes what the gate does with the files of ECG and imaging
items.

\begin{table}[htbp]
\caption{What the privacy gate does with each kind of attachment (summary).}
\label{tab:ecg-imaging-gate}
\begin{tabularx}{\linewidth}{@{}P{3.4cm}L@{}}
\toprule
{\tabhead Kind} & {\tabhead In Privacy Mode} \\
\midrule
WFDB header (\file{.hea}) & Identifying comments removed; dates moved back; ages capped; the
  record's and signal files' names replaced by tokens. \\
WFDB samples (\file{.dat}) & Sent as they are: they hold only numbers. \\
GE MUSE XML & Demographics removed except age and sex; the order removed; staff names replaced;
  every date moved back. Any other XML file is refused. \\
DICOM & The header rewritten following DICOM PS3.15 Annex E (the Basic Application Level
  Confidentiality Profile, retaining longitudinal temporal information with modified dates); the
  pixel data copied byte for byte; marked as possibly identifying. Deflated and big-endian files
  are refused. \\
EDF, EDF+, BDF, BDF+ & The header de-identified and the texts of annotations removed; the
  samples sent as they are (\cref{sec:protecting-edf}). \\
Pictures, videos, PDF & Sent as they are, marked as possibly identifying: the page shows an amber
  warning triangle on the viewer's tab. \\
Anything else & Refused: JSON ECGs, text files, and signal files with an extension other than
  \file{.dat}. \\
\bottomrule
\end{tabularx}
\end{table}

What the pixels show is not changed: ultrasound images, secondary captures and scanned documents
may carry burned-in names. File names never leave the server. In Privacy Mode an attachment's
files are named by tokens (\code{mo-f0000.hea}, \code{mo-f0001.dat}), and a real file name is not
accepted.

The databases that \cmd{muse-py index} and \cmd{dicom-py index} write hold the patients' IDs, the
key that links each row to its patient. Keep them as PHI, with the catalog's databases, and never
in a repository. They never hold names or dates of birth.
